\cotefile{55}
\scanpage{2}
% archive sheet number: 1
% lines 1
$X$ espace de Stein compact

% lines 2
\[
 A=\Gamma(X,\underline O_X)
\]

% lines 3
\deleted{$\unsure{B|{\mathbb B}}={}$ Algèbre de présentation finie}

% lines 4–7
1°) Le foncteur $M\unsure{\leadsto|\sim} M\otimes_{A_X}\underline O_X$ de la catégorie des modules de \deleted{type} \inserted{présentation} finie sur $X$ dans la catégorie des $\underline O_X$-Modules \uline{coh.} est une équivalence de catégories.

% lines 8–10
\deleted{En effet, le foncteur est exact (?) pour voir s'il est pl. fid. il suffit de prouver que $\widetilde M=0\Rightarrow M=0$. Or}

% lines 11
Soit $M$ de présentation finie,

% lines 12
\[
 A^p - A^q - M \to 0
\]

% lines 13
alors on a

% lines 14
\[
 \underline O^p - \underline O^q - \widetilde M \to 0
\]

% lines 15–16
d'où passant \unsure{au|aux|les} $\Gamma$, et notant l'exactitude (car $X$ Stein)

% lines 17
\[
 \deleted{\unsure{\Gamma(M)-|\Gamma(M)}} \quad M\xrightarrow{\sim}\Gamma(\widetilde M)
\]

% lines 18
Donc \deleted{\unsure{par propriété Stein}} \unsure{on voit|tout}

% lines 19
\[
 \deleted{\widetilde M=0\Rightarrow\Gamma(\widetilde M)=0\quad\text{i.e. }M=0}
\]

\scanpage{3}
\mspage{- 111 -}
% archive sheet number: 2
% lines 1–6
irréductible de $f^{-1}(y)$ \deleted{\illegible{2}}, de point générique $x$, est contenue dans l'adhérence d'un voisinage ouvert affine de $x$. On peut aussi remplacer $Y$ par le schéma local de $y$, donc supposer que \deleted{\illegible{1}} $Y=\mathrm{Spec}(B)$, où $B$ est un anneau local \inserted{noethérien}, d'idéal maximal $\handwritten{\underline}{m}==\handwritten{\underline}{j}_y$, et $X=\mathrm{Spec}(A)$, où $A$ est une algèbre \deleted{\illegible{2}} \inserted{noethérienne} sur $B$ ; on peut en outre (remplaçant $B$ par son image dans $A$) supposer $B\handwritten{\subset} A$.
% slanted
\margin{\handwritten{préparatifs laborieux inutiles}}

% lines 6–11
Alors il s'agit de voir que si $\handwritten{\underline}{p}$ est \deleted{\illegible{1}} un idéal premier\deleted{\unsure{m}} minimal associé à $A\handwritten{\underline}{m}$, on a $\mathrm{rg}(\handwritten{\underline}{p})\handwritten{\leq}\mathrm{rg}(\handwritten{\underline}{m})$. Or, dans l'anneau local \deleted{$\unsure{A|{}}_{\unsure{p}}$} $A_{\handwritten{\underline}{p}}$ l'idéal maximal $\handwritten{\underline}{p}A_{\handwritten{\underline}{p}}$ est engendré par $A\handwritten{\underline}{m}$, donc admet un système de $r$ générateurs, si $r$ est le rang de $\handwritten{\underline}{m}$, ce qui démontre notre assertion.

% lines 12–15
\uline{Corollaire 1 .- Soient} \inserted{$X$ \uline{et}} $Y$ \uline{\deleted{\unsure{un}} \inserted{deux} préschémas localement noethériens}, $f:X\handwritten{\to Y}$ \uline{un morphisme \deleted{\illegible{2}}}, $Y'$ \uline{une partie fermée irréductible de} $Y$, $X'$ \uline{une composante irréductible de} $f^{-1}(Y')$, \uline{dont l'image dans} $Y$ \uline{soi\handwritten{t} dense dans} $Y'$. \uline{Alors} $\mathrm{codim}_X(X')\handwritten{\leq}\mathrm{codim}_Y(Y')$.

% lines 16–18
Si $x$ et $y$ sont les points génériques de $X'$ et $Y'$ respectivement, \deleted{\illegible{1}} l'hypothèse sign\unsure{i|j}fie que $f(x)=y$, et le cor.1 n'est qu'une autre fo\handwritten{r}me de la prop.8 .

% lines 19–22
\uline{Corollaire 2 .- Soient} $X$ \uline{un préschéma \deleted{\unsure{noethérien}} \inserted{localement noethérien},} $f$ \uline{une section de} $O_X$ \uline{au-dessus de} $X$, \uline{qui n'est identiquement nulle dans aucune composante irréductible de} $X$. \uline{Soit} $Z_f$ \uline{l'ensemble des} $x\handwritten{\in}X$ \uline{tels que} $f(x)\handwritten{=0}$ \uline{Alors toute composante irréductible de} $Z_f$ \uline{est de codimension} 1 .

% lines 23–29
Il suffit de démontrer qu'il en est ainsi dans un vois\unsure{i}nage ouvert affine de tout point de $Z_f$, donc on peut supposer \deleted{\illegible{1}} $X=\mathrm{Spec}(A)$ affine d'anneau $A$ noethérien ; par hypothèse $f$ n'est \deleted{\illegible{1}} contenu dans aucun des idéaux premiers minimaux de $A$, et tous les idéaux premier\handwritten{s} minimaux associés à l'idéal principal $(f)$ sont donc de rang 1 (``Hauptidealsatz'' de Krull ;cf.Bourbaki,Alg.comm.\qquad), ce qui prouve le corollaire. Ce dernier se déduit aussi de la prop.8, en \deleted{\illegible{1}} consi-
% slanted
\margin{\handwritten{affinisation inutile, si on remarque que pour que $x\in X$ soit \deleted{le} \unsure{pt} générique d'une comp. irréd. de $Z_f$, il faut et \unsure{su.} que $f\unsure{\in|+}\underline O_x$ \deleted{\unsure{engendre}} un \deleted{\unsure{x|l'}}idéal de définition}}

\scanpage{4}
% archive sheet number: 3
% lines 1–2
Soit $N$ un autre module de prés. finie, Alors on a diagramme commutatif à lignes exactes

% lines 3–4
\diagram{\begin{tikzcd}[ampersand replacement=\&,column sep=small,row sep=small]
0\arrow[r,dash]\&\mathrm{Hom}(M,N)\arrow[r,dash]\arrow[d,"\alpha"]\&N^q\arrow[r,dash]\arrow[d,dash,"\beta"]\&N^p\arrow[d,"\gamma"]\\
\deleted{\unsure{A}}\,0\arrow[r,dash]\&\mathrm{Hom}(\widetilde M,\widetilde N)\arrow[r,dash]\&\Gamma\widetilde N^q\arrow[r,dash]\&\Gamma\widetilde N^p
\end{tikzcd}}

% lines 5–8
et comme $\beta$ et $\gamma$ sont bijectives, de $=$ $\alpha$, donc le foncteur \inserted{$M\unsure{\leadsto|\sim}\widetilde M$} est \uline{pleinement fidèle}. Prouvons qu'il est essentiellement surjectif : en effet, si un $F$

% lines 9–10
coh. sur $X$, on sait qu'il est engendré par ses sections, i.e. $\exists\,\underline O_X^{\deleted{\unsure{p}}q}\to F\to0$ \deleted{\unsure{surjectif}} \inserted{exact}

% lines 11–13
Soit $R$ le noyau, Alors il existe de $=$ $\underline O_X^p - R - 0$, d'où une suite exacte

% lines 14
\[
 \underline O_X^{p}\mathrel{\overset{u}{-}}\underline O_X^{\deleted{\unsure{p}}q} - F - 0
\]

% lines 15–16
et par suite \deleted{\unsure{si} $u$ est donné par} \deleted{et par suite} si $u$ définit. (passant aux sections)

% lines 17
\[
 A^{\deleted{\unsure{q}}p}\mathrel{\overset{u}{\to}}A^q
\]

% lines 18–19
on prend $M=\operatorname{Coker}u$ et on trouve que $F\simeq\widetilde M$.

\scanpage{5}
\mspage{- 110 -}
% archive sheet number: 4
% lines 1–2
\margin{\handwritten{abusif}}$\dim(X\handwritten{\times}_k Y)=\dim(X)+\dim(Y)$ ; \uline{en particulier, pour tout corps} $K$, \uline{extesion de} $k$, $\dim(X\handwritten{\otimes}_k K)=\dim(X)$.

% lines 3–5
Si $(U_i)$ et $(V_j)$ sont des recouvrements ouverts affines de $X$ et $Y$, $(U_i\handwritten{\times}_k V_j)$ est un recouvrement ouvert affine de $X{}_kY$, donc on peut se borner au cas où $X$ et $Y$ sont affines (n°1,cor.1 de la prop.2).

% lines 6–11
D'autre part, si \deleted{$\unsure{X}_i$} $X_i$ (resp. $Y_j$) sont les composantes irréductibles de $X$ (resp. $Y$), $X\handwritten{\times}_kY$ est réunion des $X_i\handwritten{\times}_kY_j$, qui sont ferm\handwritten{és} (\S{} 2,n°2,prop.5), donc (n°1,prop.1 (ii)) on peut supposer $X$ et $Y$ i\handwritten{r}réductibles ; \deleted{\illegible{1}} enfin, comme $(X\handwritten{\times}_kY)_{\mathrm{red}}$ et $(X_{\mathrm{red}}\handwritten{\times}_kY_{\mathrm{red}})_{\mathrm{red}}$ sont identiques (\S{} 3,n°4,cor. de la prop.11), on peut supposer $X$ e\handwritten{t} $Y$ réduits.

% lines 11–17
Soient donc $X=\mathrm{Spec}(A)$, $Y=\mathrm{Spec}(B)$, où $A$ et $B$ sont des algèbres intègres de type fini sur $k$. Si $m=\dim(X)$, $n=\dim(Y)$, il résulte du lemme de normalisation (Bourbaki,Alg.comm.\qquad) que $A$ (re\handwritten{sp.} $B$) est \deleted{\illegible{1}} un module de type fini sur un sous-anneau $A_1=\allowbreak =k\handwritten{[}a_1,\dots,a_m\handwritten{]}$ (resp. $B_1=k\handwritten{[}b_1,\dots,b_n\handwritten{]}$) où les $a_i$ (resp. $b_j$) sont algébriquement indépendants sur $k$ ; $C=A\handwritten{\otimes}_kB$ est donc un module de type fini sur $C_1=A_1\handwritten{\otimes}_kB_1$.
% slanted
\margin{\handwritten{Autre démonstration : les $\unsure{\kappa|K}(z)$ $(z\in X\times_k Y)$ sont toutes les ext. composées des ext. $\kappa(x)$ et $\kappa(y)$ $(x\in X,y\in Y)$, d'où aussitôt le résultat vu la théorie des corps}}

% lines 17–23
Pour tout idéal premier minimal $\handwritten{\underline}{p}$ \deleted{\unsure{dans}} dans $C$\handwritten{,} on voit par suite que le degré de transcendance \inserted{sur $k$} du corps des fracti\handwritten{ons} de $C/\handwritten{\underline}{p}$ est au plus $m+n$ ; \deleted{\illegible{1}} d'autre part, en vertu du th. de Cohen-Seidenberg,il existe un idéal premier $\handwritten{\underline}{q}$ de $C$ tel que $\handwritten{\underline}{q}\handwritten{\cap} C_1=\handwritten{\{}0\handwritten{\}}$, de il y a un des idéaux premiers minimaux \deleted{\unsure{dans}} dans \deleted{\illegible{1}} $C$ tel que le corps des fractions de $C/\handwritten{\underline}{p}$ soit de degré fini sur le corps des fractions de $C_1$, et par suite de degré de transcendance $m+n$ sur $k$.

% lines 24–28
\uline{Proposition 8 .- Soient} \inserted{\handwritten{$X$,}} $Y$\handwritten{,} \uline{\deleted{un} \inserted{\handwritten{deux}} préschéma\handwritten{s} localement noethérien\handwritten{s}}, $f:X\handwritten{\to}Y$ \uline{un morphisme \deleted{de type fini}. Soient} $y\handwritten{\in}Y$, \uline{et soit $x$ le point nérique d'une composante irréductible de} $f^{-1}(y)$, \uline{considéré comme préschéma sur} $\handwritten{\underline{\kappa}}(y)$ (\S{} 2,n°7,prop.17). \uline{Alors} $\operatorname{codim}_X\handwritten{\overline}{\handwritten{\{}x\handwritten{\}}}\handwritten{\leq}\operatorname{codim}_Y\handwritten{\overline}{\handwritten{\{}y\handwritten{\}}}$ (ou, ce qui revient au même, $\dim\handwritten{\underline}{O}_x\handwritten{\leq}\dim\handwritten{\underline}{O}_y$).
% slanted
\margin{\handwritten{\unsure{généraliser|Simplifier} en la formule $\dim\underline O_x\leq\dim\underline O_y+\dim\underline O_x$ (sur la fibre}}

% lines 29
La question est locale en $x$ et en $y$, puisque toute composante
% slanted
\margin{\handwritten{\deleted{\unsure{dire que la} \unsure{différence} \unsure{codim} dans $f^{-1}(y)$ \unsure{est majorée|est minorée} par la codim de \unsure{son adh.|plus un}, \unsure{ds|si} \unsure{son|non} $X$}}}
\margin{\handwritten{donner le cas d'égalité par \uline{platitude}}}

\scanpage{6}
% archive sheet number: 5
% lines 1
$\bigl[$\uline{Lemme} A \quad $X$ espace annelé, \deleted{avec $\underline O_X$ \uline{cohérent}} \inserted{on suppose}.

% lines 2
(i) $X$ quasi-compact

% lines 3
\deleted{\unsure{(ii)} $\underline O_X$ cohérent} (ii) $\underline O_X$ cohérent, les $\underline O_{X,x}$ noethériens.

% lines 4–6
\deleted{\unsure{(ii)}} \inserted{(iii)} Pour \deleted{\unsure{tt} \unsure{suite exacte stricte}} \inserted{tout} Module \deleted{cohérent} \inserted{cohérent} $\underline J$ \deleted{\unsure{de} \unsure{dans} $\underline O_X$}, \inserted{sur $X$,} on a $H^1(X,J)=0$ \deleted{dans les cas suivants :} \deleted{$\underline J$ est libre de type fini \unsure{ou|sur les} \unsure{loc} \unsure{un}}

% lines 7
Alors, posant $\deleted{\unsure{N}}A=\Gamma(X,\underline O_X)$, le foncteur

% lines 8
\[
 M\unsure{\leadsto|\sim}\widetilde M=M\otimes_{A_X}\underline O_X
\]

% lines 9–12
de la catégorie des $A$-modules de présentation finie dans la catégorie des $\underline O_X$-Modules de prés. finie (= cohérents) est une équivalence de catégories. $\bigr]$

% lines left of 1–12
% margin note written sideways
\margin{\uline{Variante} Si tt pt de $X$ a un système fond. de voisinages \deleted{\unsure{qui}} quasi-compacts, si $X$ \deleted{étant} quasi-compact, et si pour tt \deleted{\unsure{sous}} Module $F$ \unsure{Noyau|d'Noyau} ou image d'un $\underline O^p\to\underline O^q$, $H^1(X,F)=0$, alors le foncteur $M\leadsto M\otimes_{A_X}\underline O_X$ des Modules sur $A$ dans les Modules \deleted{\unsure{quasi-coh}} ayant une \deleted{\unsure{une}} présentation \deleted{finie} \deleted{\unsure{globale}} \deleted{\unsure{est une}} \inserted{globale $\underline O^{(I)}-\underline O^{\unsure{(J)|J}}$, est} une équivalence \ldots}

% lines 13–15
\uline{Corollaire 1} Le foncteur $M\leadsto\widetilde M$ est \uline{exact}. En particulier, pour tt $x\in X$, $\underline O_{X,x}$ est \uline{plat} sur $A$.

% lines 16–18
\uline{Corollaire 2} L'anneau $A$ est \textquotedblleft cohérent\textquotedblright{} : i.e. pour tout hom $A^n - A$, le noyau est de t.f.

% lines 19–20
\uline{Corollaire 3} Si $X$ est \textquotedblleft pur\textquotedblright, \deleted{\unsure{les} $\underline O_{X,x}$} alors on a correspondance biunivoque entre idéaux maximaux de \uline{type fini} de $A$, et couples $(x,\mathfrak m)$ où $x\in X$ \deleted{\unsure{fermé}} et $\mathfrak m$ idéal maximal de $\underline O_{X,x}$.

% lines 21
\uline{Remarque} On peut remplacer (ii) (iii) par

% lines 22–23
(\deleted{\unsure{i}}$\alpha$) Pour toute suite exacte $\underline O_X^p-\underline O_X^q-F\to0$, on a exactitude pour la suite des $\Gamma$

% lines 24–25
(\deleted{\unsure{ii}}$\beta$) Tout $F$ de présentation finie est engendré par ses sections \ldots

\scanpage{7}
\mspage{- 109 -}
% archive sheet number: 6
% lines 1–2
son idéal maximal, $X=\mathrm{Spec}(B)$ a deux points, $(0)$ et $\handwritten{\underline}{p}$, ce dernier étant le seul point fermé ; on a $\dim(X)=1$ mais $\dim\handwritten{\{}0\handwritten{\}}=0$.

% lines 3–6
\uline{Corollaire 2 .- Soit} $X$ \uline{un préschéma algébrique sur un corps $k$. Pou\handwritten{r} tout point} $x\handwritten{\in}X$, \uline{on a} $\dim_x(X)=\sup_i(\dim(X_i))$\handwritten{\unsure{,|$^\circ$}}, \uline{où} $X_i$ \uline{parcourt l'ense\handwritten{m}ble des composantes irréductibles de} $X$ \uline{contenant $x$. Si} $\handwritten{\{}x\handwritten{\}}$ \uline{est ferm\handwritten{é} on a} $\dim_x X=\operatorname{codim}_X x=\dim(\handwritten{\underline}{O}_x)$\handwritten{,} \inserted{\handwritten{et en tous cas $\dim_x X=\dim\underline O_{X,x}+\dim k(x)/k$}}

% lines 7–13
En effe\handwritten{t}, \deleted{pour tout ouvert $U$ contenant \deleted{\illegible{1}} $x$}, on a (n°1,cor.de la prop.1) $\dim_x(X)=\sup_i(\dim_x(X_i))$ et \deleted{\illegible{1}} d'après le cor.1 ; $\dim_x(X_i)=\allowbreak=\dim(X_i)$. D'autre part, comme \inserted{si $\handwritten{\{}x\handwritten{\}}$ est fermé,} toute chaîne de parties fermées irréductibles contenant $x$ est contenue dans\handwritten{\,}l'un des $X_i$, on a $\operatorname{codim}_X(x)\handwritten{=}\allowbreak=\sup_i(\operatorname{codim}_{X_i}(x))$, et en vertu de la formule (6) (applicable puisque $X_i$ est strictement équidimensionnel), $\operatorname{codim}_{X_i}(x)=\dim(X_i)$, d'où \deleted{\illegible{1}} la seconde assertion du corollaire.
\margin{\handwritten{$=\qquad?$}}

% lines 14–15
\uline{Corollaire 3 .- Soient} $X,Y$ \uline{deux préschémas algébriques sur un corps $k$. Pour tout $k$-morphisme} $f:X\handwritten{\to}Y$, \uline{on a} $\dim(\handwritten{\overline}{f(X)})\handwritten{\leq}\dim(X)$.
% slanted
\margin{\handwritten{\unsure{déjà vu en|a-t-on} \S{}1}}
\margin{\handwritten{\uline{égalité dans le cas quasi fini}}}

% lines 16–23
En considérant les composantes irréductibles de $\handwritten{\overline}{f(X)}$, on est aussitôt ramené au cas où \deleted{\illegible{2}} $Y$ est irréductible et $f$ dominant (n°\handwritten{1,} prop.1 (ii)) ; en outre, une au moins des composantes irréductible\handwritten{s} $X_i$ de $X$ est telle \inserted{alors} que \deleted{\illegible{1}} $f(X_i)$ soit dense \deleted{\illegible{1}} dans $Y$, donc on peu\handwritten{t} aussi supposer $X$ irréductible. Alors, si $x$ est le point générique de $X$, $f(x)=y$ est le point générique de $Y$, et comme il \deleted{\illegible{1}} correspond à $f$ un $k$-monomorphisme de $\handwritten{\underline{\kappa}}(y)$ dans $\handwritten{\underline{\kappa}}(x)$, on a \deleted{\illegible{1}} $\mathrm{deg.tr.}_k\,\handwritten{\underline{\kappa}}(y)\handwritten{\leq}\mathrm{deg.tr.}_k\,\handwritten{\underline{\kappa}}(x)$, d'où le corollaire.

% lines 24–29
Ici encore, ce résultat ne s'étend pas à des schémas intègres que\handwritten{l}conques. Par exemple, soient $B$ un anneau de valuation discrète, $K$ son corps des fractions, $X=\mathrm{Spec}(K)$, $Y=\mathrm{Spec}(B)$ ; à l'injection canonique $B\quad K$ correspond u\unsure{x|n} morphisme $f:X\handwritten{\to}Y$ \uline{de type fini} ($K$ étant en effet engendré par $1/\handwritten{\pi}$, si \deleted{\illegible{1}} $(\handwritten{\pi})$ est l'idéal maximal de $B$) ; mais on a \deleted{\illegible{1}} $\dim(Y)=1$ et $\dim(X)=0$.
% slanted
\margin{\handwritten{\deleted{\uline{Corollaire} $\dim_x X=\dim\underline O_x+\dim k(x)/k$}}}
\margin{\handwritten{déjà vu en \S{}1}}

% lines 30
\uline{Corollaire 4 .- \unsure{S}oient} $X,Y$ \uline{deux préschémas algébriques sur $k$ ; \handwritten{on a}}

\scanpage{8}
% archive sheet number: 7
% lines 1–2
\unsure{On|[On|Du} trouve en tous cas, sous la seule hyp. ($\alpha$) de la remarque, les conclusions :

% lines 3–6
$B\unsure{\leadsto|\sim} B\otimes_{A_X}\underline O_X$ est une équivalence de la catégorie des $A$-Alg. de présentation finie avec la catégorie des $\underline O_X$-Alg. de \uline{présentation} \uline{globale} finie.

% lines 7–11
Si $B$ \unsure{et|est} \deleted{\unsure{une}} $\underline B$ se correspondent, Alors $M\leadsto M\otimes_{B_X}\underline B$ est une équiv. de la catégorie des $B$-Modules de présentation finie, et la catégorie des $\underline B$-Modules des présentation \uline{globale} finie.

% lines 12–13
\uline{Question} Soit $\underline B$ de \inserted{Alg.} prés. finie sur $X$, le noyau $J$ de $\underline O_X\unsure{-|\to}B$ est-il cohérent ??

% lines 14–18
C'est une question de pure Algèbre comm., peut s'énoncer : si $A$ est un anneau cohérent, \unsure{et|\deleted{$A$}} $J$ un idéal de type fini dans $B=A[T]$, sa trace $I$ sur $A$ est-elle de type fini ??

% lines 18–20
[ Plus gén, est-il vrai que les $I_{\unsure{m|n}}\subset B_m$ sont des $A$-modules de type fini ?? ]

% lines bottom right
\margin{OK. ds le \textquotedblleft cas classique\textquotedblright{} $X$ \deleted{\unsure{un}} \unsure{polyèdres|polyèdre|polycylindre} \unsure{analytiques|analytique} \ldots}

\scanpage{9}
% archive sheet number: 8
% lines 1
\[
 \operatorname{Ass}F=\handwritten{\bigcup_i}\operatorname{Ass}F_i
\]

% lines 2
(iii) $\operatorname{Ass}\handwritten{\coprod} F_i=\handwritten{\bigcup_i\operatorname{Ass}}F_i$ \qquad pour toute famille $(F_i)$ de modules qu c

% lines 3
Beweis klaar.

% lines 4–6
\uline{Proposition} 6.7. \quad Supposons $X$ loc noeth. Pour que la section $f\varepsilon I\underline O_X$ definisse une homothetie $F\quad F$ qui est \uline{injective}, il f et s que $X_f\quad\operatorname{Ass}F$, i.e. que $V(f)\quad\operatorname{Ass}F=\varnothing$.

% lines 7–8
En effet, on est ramene au cas affine (par exemple) o\handwritten{\`}{u} c'est dans Bourbaki. On en conclut:
% slanted
\margin{\handwritten{démonstrat\unsure{ion} directe}}

% lines 9–11
\uline{Corollaire} 6.8. \quad Soit $I$ un Ideal coherent sur $X$ loc noeth, $F$ un Module coherent sur $X$. Pour que $\underline I$ soit 'non diviseur de 0 dans $F$' il f et s que $V(I)\quad\operatorname{Ass}\unsure{\mathbb F|F}=\varnothing$.

% lines 12–13
On peut mettre ici si on veut IV 2.8.18. et 2.8.19. (bien qu'il soit possible qu'on ne les utilisera plus par la suite ...)

% lines 14–17
\uline{Proposition} 6.9. \quad Soient $X$ preschema loc noeth, $U$ un ouvert dense. Pour que $X$ soit reduit, il f et s que $U$ le soit, et que $X$ soit sans 'cycles premiers immerges'$^*$, ou encore que $X$ soit sans cycles premiers imm. et que l'on ait \quad ait long $\underline O_x=1$ pour $x$ point gen d'une comp irr de $X$.

% lines 18
(Cette notion * devrait etre definie, pour un Module, dans 6.1.).

% lines 19–26
Evidemment (i) \quad (ii) \quad (iii), d'autre part (iii) \quad (i) \deleted{puisque \unsure{d'}} d'apres le cas affine, bien connu \deleted{\unsure{(0} etant intersection des ideaux primaires \unsure{asso}} ou en notant qu'une sectionndu faisceau $\underline N$ des ideaux nilpotents a unnsupport qui ne peut contenir aucun point generique, donc s'il n'etait vide definirait par ses composantes irreductibles des cycles associes immerges. Il est peut-etre preferable d'essayer de donner \deleted{plutot} de telles demonstrations 'intuitives', plutot que de se reduire toujours a l'enonce connu dans le cas affine.

% lines 27–28
Nous allons moduler la notion de preschema reduit, en disant qu'un \uline{Module} qu coh $F$ \uline{est reduit} s'il est \uline{de type fini}, \uline{sans cycle premier associe}

\scanpage{10}
% archive sheet number: 9
% lines 1–4
déjà que le foncteur $M\unsure{\leadsto|\sim|\mapsto}\widetilde M$ des $\Gamma B$-\deleted{\unsure{M|$\mathcal M$}}modules de présentation finie dans les $B$-Modules de prés. finie est pl. fidèle.

% lines 4–8
\deleted{\unsure{Pour voir|Prouvons} Reste à prouver que si $M$ est un $B$-Module de prés. finie, il est engendré par ses sections. Soit $x$ un pt fermé de $X$ [\inserted{N.B.} (Nous supposons $X$ \textquotedblleft pur\textquotedblright{} ! ]. Alors $M_x$ est un $B_x$-Module de présentation finie}

% lines 9
\deleted{Soit $M'$ le sous-module engendré par les sections.}

% lines 10–14
(en utilisant \unsure{seult} la condition $\alpha$) de la \textquotedblleft remarque\textquotedblright{} ci-dessus). Reste à prouver $\beta$) [qui \unsure{en conclut|ne cloche|est inclus} dans le cas $\underline O_X$ \unsure{fixé|fini} sur \unsure{$X$|$\hat x$} compact \ldots !.], sous les conditions plus fortes \unsure{précisées|possibles}.

% lines 14–17
Il faut prouver que $B$ satisfait conditions (ii) et (iii), pour (ii) c'est l'hyp, puis (iii)$'$ c'est encore un passage à la limite, car

% lines 18
\[
 J=\varinjlim_m J_m
\]

% lines 19
il faut prouver que les $J_m\inserted{=J\cap B_m}$ sont cohérents
% slanted
\margin{? histoire d'anneaux cohérents ??}

% lines 20
-----

\scanpage{11}
\mspage{- 101 -}
% archive sheet number: 10
% lines 1
$X$\handwritten{\,}\uline{des composantes irréductibles de} $Y$.
% slanted
\margin{\handwritten{codim. en \uline{un} point ...}}

% lines 2–3
\uline{Proposition 1} .- (i) \uline{Si} $Y$ \uline{est une partie d'un espace topologique} $X$ $\dim(Y)\handwritten{\leq}\dim(X)$.

% lines 4–5
(ii) \uline{Si un espace topologique} $X$ \uline{est réunion finie de parties fermé\handwritten{es}} $X_i$, \uline{alors} $\dim(X)=\sup_i(\dim(X_i))$.

% lines 6
\deleted{(iii) \uline{Pour toute \unsure{application continue}} \illegible{4} $\unsure{f(X)}$ \unsure{on a} \unsure{$(i)$}.}

% lines 7–10
Pour toute partie irréductible $Z$ fermée dans $Y$, l'adhérence $\overline Z$ de $Z$ dans $X$ est irréductible ; d'où (i). Si $X=\handwritten{\bigcup_i}X_i$, les $X_i$ étant f\handwritten{er}més, toute chaîne de parties fermées irréductibles de $X$ est nécessairement contenue dans l'un des $X_i$, ce qui prouve (ii).

% lines 10–12
\deleted{\unsure{Pour} (iii) \illegible{7} $\unsure{Y_x}$ \unsure{est une} \illegible{3} \unsure{application de} $\unsure{X}$ \unsure{(i)}.}

% lines 12–17
On notera que si $f$ est une application continue de $X$ dans un espace topologique $Y$, on peut avoir $\dim(f(X))\handwritten{>}\dim(X)$ ; on en a un exemple en prenant pour $X$ un espace discret à 2 éléments $a,b$, pour $Y$ l'ensemble $\handwritten{\{}a,b\handwritten{\}}$ muni de la topo\handwritten{lo}gie dans laquelle $\varnothing$, $\handwritten{\{}a\handwritten{\}}$ et $\handwritten{\{}a,b\handwritten{\}}$ sont les seuls ensembles fermé\handwritten{s} si $f$ est l'application identique $X\handwritten{\to}Y$, $\dim Y=1$ et $\dim X=0$.

% lines 18–20
\uline{Corollaire .- Soient} $x\handwritten{\in}X$, $U$ \uline{un voisinage de} $\handwritten{x}$, $Y_i$ $(1\handwritten{\leq}i\handwritten{\leq}n)$ \uline{des parties fermées de} $U$ \uline{telles que} $x\handwritten{\in}Y_i$ \uline{pour} $1\handwritten{\leq}i\handwritten{\leq}n$ \uline{et que} $U$ \uline{soit l\handwritten{a} réunion des} $Y_i$. \uline{Alors} $\dim_x X=\sup_i(\dim_x(Y_i))$.

% lines 21–25
En effet, d'après la prop.1, $\dim_x(X)=\inf_V(\sup_i(\dim(Y_i\handwritten{\cap}V)))$ lorsqu\handwritten{e} $V$ parcourt l'ensemble des voisinages ouverts de $x$ contenus dans $U$. Le corollaire est évident si $\sup_i(\dim_x(Y_i))=+\handwritten{\infty}$ ; sinon, il y a un voisinage $V_0$ de $x$ tel que \deleted{$\unsure{\sup_i|X_i}$\illegible{1}} $\dim(Y_i\handwritten{\cap}V)=\dim_x(Y_i)$ pour tout $\handwritten{i}$ et tout $V\handwritten{\subset}V_0$.

% lines 26–27
\uline{Proposition 2} .- (i) \uline{Soit} $\handwritten{\Phi}$ \uline{l'ensemble des parties fermées irrédu\handwritten{cti}bles de} $X$. \uline{On a} $\dim(X)=\sup_{Y\handwritten{\in\phi}}(\operatorname{codim}_X(Y))$.

% lines 28
(ii) \uline{On a} $\dim(X)=\sup_{x\handwritten{\in}X}\dim_x(X)$.

% lines 29
(i) découle aussitôt des déf.1 et 2 . \deleted{\unsure{D'autre part il existe dans}}

\scanpage{12}
% archive sheet number: 11
% lines 1–2
\deleted{\uline{Corollaire}} \deleted{\uline{Proposition}} \inserted{\uline{Corollaire}} Supposons que l'on ait la condition \deleted{\unsure{(ii)}} suivante plus forte que (ii)

% lines 3–4
(ii$'$) \inserted{Les $\underline O_{X,x}$ noeth, et} Pour tt entier $n$, $\underline O_X[T_1,-T_n]$ est un faisceau d'anneaux cohérent.
\margin{?}

% lines 5
\inserted{(iii$'$) Tout $x\in X$ a un syst. fond. de voisinages \unsure{qu.-cpts|qui-compts} satisfont la condition (iii) pour le faisceau induit par $\underline O_X$ ....}

% lines 6–9
Alors la conclusion du lemme A est valable en y remplaçant $\underline O_X$ par $\underline O_X[T_1,-T_n]$, plus gén. par n'importe \unsure{quelle|quel} faisceau d'Algèbres $\underline B$ de présentation finie.

% lines 10–13
Soit $\underline B^p\mathrel{\overset{\inserted{u=(u_{ij})_{(i,j)\in P\times Q}}}{-}}B^q-M-0$ une suite exacte, \inserted{$\Sigma$} \deleted{\unsure{Alors on}} de $B$-Modules ($B=\underline O_X[T_1,-,T_n]$) Alors on a pour tt entier $m$ une suite exacte \inserted{$\Sigma_m$} de $\underline O_X$-Modules \inserted{cohérents}

% lines 14
\[
 \underline B_m^p - B_{N+m}^q\to M_m-0
\]

% lines 15
($N$ majorant des degrés des $u_{ij}$)

% lines 16
\deleted{\unsure{alors}} et $\Sigma$ est \uline{lim} des $\Sigma_m$. On aura

% lines 17
\[
 \Gamma B_m^p - \Gamma B_{\unsure{m+m|N+m|m+N}}^q - \Gamma M_m\to0
\]

% lines 18
d'où en passant à la limite

% lines 19
\[
 \Gamma B^p - \Gamma B^q - \Gamma M-0
\]

% lines 20
Ce qui implique $\alpha$. \deleted{\unsure{P}} Ceci prouve

\scanpage{13}
\mspage{- 102 -}
% archive sheet number: 12
% lines 1–4
D'autre part, si $Z_{\unsure{o|0}}\handwritten{\subset} Z_1\handwritten{\subset}\ldots\handwritten{\subset}Z_n$ est une chaîne de parties fermées irréductibles de $X$ et si $x\handwritten{\in}Z_o$, pour tout \unsure{o|O}uvert $U\handwritten{\subset}X$ contenant $x$, $U\handwritten{\cap}Z_i$ est irréductible et fermé dans $U$, et comme $\overline{U\handwritten{\cap}Z}_i=Z_i$, les $U\handwritten{\cap}Z_i$ sont distincts ; donc $\dim(U)\handwritten{\geq}n$, ce qui démontre (ii).

% lines 5–6
\uline{Corollaire 1 .- Si} $(U_{\handwritten{\alpha}})$ \uline{est un recouvrement ouvert de} $X$, \uline{on a} $\dim(\handwritten{X)}=\sup_{\handwritten{\alpha}}(\dim(U_{\handwritten{\alpha}}))$.

% lines 7–8
\uline{Corollaire 2 .- Soit} $X$ \uline{un espace noethérien vérifiant l'axiome} $(T_o)$ \uline{et soit} $F$ \uline{l'ensemble des points fermés de $X$.Alors} $\dim(X)=\sup_{x\handwritten{\in}F}(\dim_x(X)\handwritten{)}$

% lines 9–10
Avec les notations de la démonstration de la prop.2 (ii), il suffi\handwritten{t} d'observer qu'il existe dans $Z_o$ un point fermé (\S{} 4,n°1,lemme 1).

% lines 11
\uline{Corollaire 3 .- Soit} $Y$ \uline{une\deleted{\unsure{s}} partie fermée de} $X$. \uline{Alors}

% lines 12
\[
 \dim(Y)+\operatorname{codim}_X(Y)\handwritten{\leq}\dim(X).\tag*{(1)}
\]

% lines 13–16
Toute chaîne de parties fermées irréductibles de $Y$ étant contenue dans une composante irréductible de $Y$, on est ramené (par la déf.2) au cas où $Y$ est irréductible, et (1) est alors conséquence des définitions 1 et 2 .

% lines 17–18
\uline{Corollaire 4 .- Soient} $Y,Z,T$ \uline{trois parties fermées irréductibles de} $X$ \uline{telles que} $Y\handwritten{\subset}Z\handwritten{\subset}T$. \uline{Alors}

% lines 19
\[
 \operatorname{codim}_T(Y)\handwritten{\geq}\operatorname{codim}_Z(Y)+\operatorname{codim}_T(Z).\tag*{(2)}
\]

% lines 20
C'est une conséquence immédiate de la déf.2 .

% lines 21–24
\uline{Proposition 3 .- Soit} $X$ \uline{un espace noethérien vérifiant l'axiome} $(T_o)$ \uline{Pour que} $\dim(X)=0$, \uline{il faut et il suffit que} $X$ \uline{soit fini et discret}. \uline{Pour qu'une partie fermée} $Y$ \uline{de} $X$ \uline{soit de codimension nulle dans} $X$, \uline{il faut et il suffit que} $Y$ \uline{contienne une composante irréductible de}

% lines 25–29
La seconde assertion est immédiate (et indépendante de toute hypothèse sur $X$). De même, si $X$ est séparé (et a fortiori discret), il est clair que $\dim(X)=0$, les seuls ensembles fermés irréductibles étant les points. Inversement, si $X$ est noethérien et satisfait à $(T_o)$, toute composante irréductible de $X$ contie\unsure{n}t un point fermé \inserted{\handwritten{(\S{} 4, n°1, lemme 1)}},

\scanpage{14}
% archive sheet number: 13
% lines 1–2
\deleted{\unsure{Prop}} \uline{Lemme B} Supposons \deleted{de plus} que $X$ satisfait les conditions suivantes : (i) (ii) \uline{et}

% lines 3
(iv) Les $\underline O_{X,x}$ $(x\in X)$ sont \uline{noethériens}

% lines 4–10
(v) Pour tt \deleted{\unsure{sous-module}} Idéal coh. $J$, \deleted{\unsure{tout point}} posant $Y=(\operatorname{Supp}\underline O_X/J,\underline O_X/J)$, et tt $y\in Y$ tel que $\underline O_{Y,y}$ est \uline{intègre}, il existe un voisinage ouvert $U$ de $y$ \inserted{dans $Y$} ayant la propriété suivante : toute section $f$ de $\underline O_Y$ sur \deleted{un voisinage} $U$ qui est \deleted{telle que} nulle au voisinage de $y$, est nulle sur $U$.

% lines 10–13
Alors pour toute \deleted{suite} $F$ coh. sur $X$, \deleted{\unsure{et}} tte suite \inserted{croissante} $F_i$ de sous-faisceaux coh. de $F$ est stationnaire \ldots

% lines 14–15
\uline{Corollaire} Sous les conditions (i) à (v), $A=\Gamma(X,\underline O_X)$ est un anneau \uline{noethérien}.

% lines 16–18
Ces conditions sont satisfaites si $X$ est un \deleted{espace de Stein} \inserted{\unsure{polyèdre|polyèdres|polycylindre} analytique} (compact) \deleted{\unsure{dont on trouve}} sur $\mathbb C$. \deleted{dont} t\unsure{t} \deleted{\textquotedblleft \unsure{presque|faisceau} \unsure{tout|étalé} fermé\textquotedblright{} est loc connexe} \ldots
% slanted
\margin{préciser ce point}

\scanpage{15}
\mspage{- 104 -}
% archive sheet number: 14
% lines 1–2
$\operatorname{codim}_{Z_1}(Z_{\unsure{o|0}})=1$, il résulte de b) que $\operatorname{codim}_{Z_n}(Z_1)=\operatorname{codim}_{Z_n}(Z_o)-1\handwritten{>}n-1$ ce qui contredit l'hypothèse de récurrence.

% lines 3–4
Lorsque les conditions de la prop.4 sont remplies, on dit que $X$ stisfait à la \uline{condition des chaînes}.

% lines 5–6
\uline{Proposition 5 .- Soit} $X$ \uline{un espace noethé\unsure{r}ien, de dimension finie, vérifiant l'axiome} $(T_o)$. \uline{Les conditions suivantes sont équivalente\handwritten{s}}.

% lines 7–8
a) \uline{Deux chaînes maximales de parties fermées irréductibles de} $X$ \uline{on\handwritten{t} même longueur}.

% lines 9–11
b) \uline{\deleted{\unsure{Toutes les composantes irréductibles} \illegible{4}} \inserted{$X$ est équidimensionnel \deleted{\handwritten{et}} \handwritten{et équicodimensionnel}}, \deleted{les poin\handwritten{ts} fermés de $X$ ont même codimension}, et} $X$ \uline{satisfait à la condition des chaînes}.

% lines 12–13
c) \deleted{\unsure{Les composantes irréductibles de} \illegible{3} \unsure{de} $X$} $X$ \uline{est équidimensionnel et pour deux parties fermées irréductibles} $Y\handwritten{\subset}Z$ \uline{de} $X$, \uline{on a}

% lines 14
\[
 \dim(Z)=\dim(Y)+\operatorname{codim}_Z(Y).\tag*{(4)}
\]

% lines 15–16
d) \inserted{\handwritten{$X$ est équicodimensionnel}} \uline{\deleted{Les codimensions des points fermés de $X$ sont les mêmes}, et pou\handwritten{r} deux parties fermées irréductibles} $Y\handwritten{\subset}Z$ \uline{de} $X$, \uline{on a}

% lines 17
\[
 \operatorname{codim}_X(Y)=\operatorname{codim}_Z(Y)+\operatorname{codim}_X(Z).\tag*{(5)}
\]

% lines left of 17–24
% slanted
\margin{\handwritten{\uline{\unsure{Dire|Dim}} \uuline{Cor}. Soit $X_i$ une comp. irréd. de $X$, et $x$ un pt fermé de $X$. On a $\dim X=\operatorname{codim}_X x=\dim X_i$.}}

% lines 18–25
Les hypothèses sur $X$ entraîne\unsure{n}t que les extrémités d'une chaîne maximale de parties fermées irréductibles sont nécessairement une composante irréductible de $X$ et un point fermé ; en outre, toute chaîne saturée \deleted{\illegible{1}} \inserted{d'extrémités} $Y$ et $Z$ \inserted{($Y\handwritten{\subset}Z$)} est contenue dans une chaîne maximale dont les éléments \inserted{autres que ceux de la chaîne donnée} sont ou bien contenus dans $Y$ ou bien contiennent $Z$. Ces remarques établissent aussitôt l'équivalence de a) et b),et prouvent aussi que \inserted{si a) est vérifiée,} pour toute partie fermée irréductible $Y$ de $X$, o\handwritten{n} a

% lines 26
\[
 \tikz[baseline=(eq.base)]{
  \node[inner xsep=4pt,inner ysep=2pt] (eq) {$\dim(Y)+\operatorname{codim}_X(Y)=\dim(X),$};
  \draw[thin] plot[smooth cycle] coordinates {(eq.west) (eq.north west) (eq.north) (eq.north east) (eq.east) (eq.south east) (eq.south) (eq.south west)};
  \draw[thin] ([xshift=-28mm,yshift=1mm]eq.west) to[bend left=3] (eq.west);
}\tag*{(6)}
\]

% lines 27–29
et comme (3) a aussi lieu, (\deleted{\unsure{5}}\handwritten{4}) et (\deleted{\unsure{6}}\handwritten{5}) s'en déduisent aussitôt. Inversement, (4) entraîne (3), donc la condition des chaînes (prop.\handwritten{4)}) et en outre, en appliquant (4) pour \deleted{\illegible{1}} $Y$ \deleted{\illegible{1}} réduit à un point

\scanpage{17}
% archive sheet number: 16
% lines 1–4
Il existe un tore complexe, de dim (complexe) $2$, fibré sur un tore complexe de dimension $1$, \deleted{\unsure{à fibres|en fibres} \illegible{1}} par un sous-tore complexe de dim $1$, qui n'est pas une variété abélienne.

% lines 5–7
(À fortiori, un fibré analytique à base et fibre des variétés de Hodge, n'est pas nécessairement une variété de Hodge)

% lines 8–9
Exemple : On prend le tore $T$ quotient de $C^2$ par le groupe $\Gamma$ engendré par

% lines 10
\[ e_1,\,ie_1,\,e_2,\,e_1+ike_2 \]

% lines 11–13
avec $k$ \uline{irrationnel}, et le sous-tore $T_0$ image de l'espace complexe $\unsure{V|V'}$ engendré par $e_1$, (donc $T_0=C/Z^2$).

% lines 13–17
Si $T$ était abélienne, on sait qu'il existerait un supplémentaire $C$-vectoriel $V'$ de $\unsure{V|V'}$, tel que $V$ soit $\Gamma$-propre, i.e. $V$ engendré par $\Gamma\cap V$. À fortiori, il

% lines 18–21
existerait $x\in\unsure{\Gamma\cap V|V'|V|\deleted{\Gamma\cap V}V'}$ \deleted{tel que $\unsure{ix|\in}$} \deleted{$ix$ soit de la forme \deleted{$\unsure{\rho x+\sigma y|\rho x+\sigma ux}$, \illegible{2}} \inserted{$\rho x+\sigma ux$} $\rho$ et $\sigma$ réels, \deleted{$y\in\Gamma\cap V$} $u$ complexe, $ux\in\Gamma$, (prendre $x\in\Gamma\cap V$ non nul)}. \unsure{Soit|Si|Le} $u=u'+iu''$

\scanpage{18}
% archive sheet number: 17
% lines 1–3
et un nb complexe \uline{non réel} $c=c'+ic''$ $(c',c''\in\unsure{R|\mathbb{R}})$ tel que $x$ et $cx\in\Gamma$. \inserted{(prendre $x\in V'\cap\Gamma$, $y\in V'\cap\Gamma$ \deleted{\unsure{si|soit} \illegible{1}} lin. ind. sur $R$)}

% lines 3
Les coordonnées de $x$ sont

% lines 4–5
\[ \alpha+\delta,\quad\beta,\quad\gamma,\quad k\delta,\qquad
\begin{aligned} &(\gamma\text{ ou }\delta\ne0\deleted{\unsure{,|)|;}}\\ &\alpha\beta\gamma\delta\in Z) \end{aligned} \]

% lines 6
celles \deleted{$\unsure{\beta|\gamma|\substack{\beta\\\gamma\\\delta}}$} de $cx$ sont

% lines 7–8
\[ \begin{aligned} &c'(\alpha+\delta)\unsure{-|+}c''\beta,\quad c''(\alpha+\delta)+c'\beta,\\
&c'\gamma-k\delta c'',\quad c''\gamma+c'k\delta \end{aligned} \]

% lines 9–14
Il faudrait donc que ce dernier vecteur puisse s'écrire $\alpha'+\delta',\beta',\gamma',k\delta',$ avec $\alpha'\beta'\gamma'\delta'\in Z$. En particulier, les premières $2$ composantes devraient être entiers, d'où on conclut que $c'$ et $c''$ seraient \uline{rationnels}.

% lines 14–16
Pour que \unsure{la|les} \deleted{deux} \deleted{deuxième} \inserted{troisième} \unsure{composante|composantes} \unsure{puisse|puissent} être \unsure{rationnelle|rationnelles}, il faut donc $\deleted{\unsure{k|\gamma}}\delta c''=0$ \deleted{\unsure{ceci|c'est} \illegible{1}}

% lines 17–18
Comme $c''\ne0$, cela implique $\delta=0$. Donc les composantes de $cx$ sont

% lines 19
\[ \alpha c'-\beta c'',\quad\beta c'+\alpha c'',\quad c'\gamma,\quad c''\gamma \]

% lines 20–23
On en tire $k\delta'=c''\gamma$, ce qui n'est possible que si les deux mbres sont nuls, d'où $\gamma=0$. Mais, \deleted{\unsure{donc}} \inserted{comme $\gamma=\delta=0$} $x$ serait dans \unsure{l'espace|l'esp.} engendré par $\unsure{e_0|e_1}$ et $ie_1$, absurde
